\chapter{Методы}

\section{Спектральный анализ}

Базовый протокол эксперимента был сначала реализован для ResNet18, затем адаптирован для ViT и Swin. Обобщённая процедура вычисления спектральной характеристики архитектуры приведена в алгоритме~\ref{alg:architecture_spectral}.

\begin{algorithm}[h]
\caption{Обобщённая процедура вычисления спектральной характеристики архитектуры}
\label{alg:architecture_spectral}
\vspace{1em}
\begin{algorithmic}[1]
\Require Выбранный слой нейронной сети, геометрические параметры $H, W$.
\Ensure Нормализованная радиальная кривая спектра мощности $\tilde E(\rho)$.
\vspace{1em}

\State Получить представление $\mathcal{F}$ из выбранного слоя/блока
\If{$\mathcal{F}$ является 4D тензором (CNN)}
    \State Оставить тензор без изменений: $(B,C,H,W)$
\Else{ \ $\mathcal{F}$ является 3D тензором вида $(B, N, D)$ (ViT, Swin)}
    \State Выполнить изменение формы в $(B, H, W, D)$ с $H=W=\sqrt{N}$
\EndIf
\State Центрировать тензор по пространству
\State Выполнить 2D FFT по осям $H,W$ для каждого канала и сдвиг центра 
\State Вычислить спектр мощности: $S(b, u, v, c) \gets |FFT(\dots)|^2$
\State Усреднить спектр по каналам: $\bar{S}_b(u, v)$
\State Для каждого радиального кольца $A_j$ вычислить среднюю мощность $M_j$ для визуализации
\State Для каждого радиального кольца $A_j$ вычислить суммарную энергию $E_j$ для метрик
\State Нормализовать суммарную энергию: $p_j = E_j / \sum_j E_j$
\State \Return $\tilde E(\rho)$
\end{algorithmic}
\end{algorithm}

Первым этапом является процедура извлечения представлений. В данной работе «представление» — это набор значений в некотором слое модели, организованный в пространственную структуру. Спектральный анализ опирается на предположение, что эти активации можно интерпретировать как дискретную функцию на двумерной решётке.

Для каждой из исследуемых моделей были реализованы функции перехвата (hook) на выходах ключевых блоков: \texttt{layer1}--\texttt{layer4} для ResNet и равномерно по глубине для трансформеров ViT и Swin. На каждый такой блок подаётся один и тот же набор изображений из валидационной выборки ImageNet \cite{priyerana_imagenet10k} (после стандартной нормализации). 
Реализация функций перехвата отличается в зависимости от используемой архитектуры, поскольку различаются представления, с которыми оперируют модели. Так, например, в ResNet каждый слой явно оперирует с двумерной картой признаков: тензоры имеют размерность $(B, C, H, W)$, где $H$ и $W$ — пространственные размеры карты признаков на данном слое, $C$ — число каналов, $B$ — размер батча. 

ViT работает иначе: изображение разбивается на непересекающиеся патчи размера $16\times16$. Каждый патч линеаризуется в вектор и пропускается через линейный слой, превращаясь в токен размерности $D$. В результате получается последовательность токенов, длина которой составляет $N = W_{patch}\times H_{patch}$, где $W_{patch}, H_{patch}$ - количество патчей по горизонтали и вертикали. Внутри самой модели перед первым блоком к токенам добавляются позиционные эмбеддинги; также используется CLS-токен, поэтому на вход поступает $N+1$ токенов. Токены обрабатываются механизмом самовнимания, в котором пространственная близость не задаётся так же явно, как в свёрточной карте признаков. Поэтому на выходе блока получается не двумерная карта, а одномерная последовательность размерности $(B, N+1,D)$ с фиксированным порядком токенов. 

Чтобы провести для ViT спектральный анализ, сопоставимый с анализом CNN, в функции перехвата последовательность токенов преобразуется в двумерную карту признаков. В классификационных вариантах ViT дополнительно используется CLS-токен. Он не соответствует отдельной пространственной области изображения и не является позиционным эмбеддингом; его задача состоит в агрегировании информации для последующего предсказания класса.
Поэтому для проведения спектрального анализа CLS-токен исключается, так как он не имеет координаты на двумерной решётке патчей. Размеры решётки патчей определяются как $H_{patch} = W_{patch} = \sqrt{N}$, поскольку после стандартной нормализации все изображения имеют квадратный размер $224\times 224$. Затем тензор преобразуется из формы $(B, N, D)$ в $(B, C, H_{patch}, W_{patch})$ построчной развёрткой таким образом, что $i$-й токен соответствует позиции $(i\mod W_p ,\lfloor i / W_p \rfloor)$. Важно учитывать, что такая карта признаков не эквивалентна карте признаков CNN. Пространственное разрешение в ViT равно количеству патчей, а не пикселей, поскольку каждый токен соответствует неперекрывающемуся блоку пикселей и представляет его агрегированное векторное описание без явного сохранения внутренней структуры патча.

В архитектуре Swin входное изображение разбивается на непересекающиеся патчи размера $4\times4$. Как и в CNN, по мере углубления модели уменьшается пространственное разрешение: между стадиями выполняется слияние патчей, которое увеличивает количество каналов и уменьшает пространственное разрешение вдвое. В Swin отсутствует CLS-токен, поэтому преобразование выполняется прямым восстановлением решётки размером $\sqrt{N} \times \sqrt{N}$. Следует отметить, что механизм сдвига окон для спектрального анализа не требует специальной дополнительной обработки, так как он уже учтён внутри блока.

Перед применением каких-либо частотных преобразований необходима предобработка.
Ключевым шагом является вычитание пространственного среднего из каждой карты признаков:

\begin{equation}
    \tilde{x}_{b,c}(h,w) = x_{b,c}(h,w) - \frac{1}{HW} \sum_{h,w} x_{b,c}(h,w)
\end{equation}

Эта операция необходима для устранения нулевой частоты (постоянной составляющей, или DC-компоненты) из сигнала. DC-компонента представляет собой среднее значение по пространству и по абсолютной величине часто на порядки превосходит энергию остальных частот, тем самым маскируя вклад более высоких частот. Кроме того, постоянная составляющая несёт информацию, которая для целей данного исследования является скорее помехой. Её удаление необходимо, чтобы анализ был чувствительным именно к вариациям признаков, а не к их абсолютному уровню.
В рамках данной работы операция выполняется независимо для каждого канала и каждого образца в батче: из каждого элемента двумерной матрицы вычитается среднее арифметическое по всем элементам этой же матрицы. Альтернативой могло бы быть высокочастотное фильтрование, но оно требует подбора граничных частот и вносит дополнительный произвол; вычитание же среднего — это простейший и в то же время строгий способ обнуления DC-компоненты без искажения остального спектра.

Для перехода из пространственной области карт признаков в частотную используется двумерное дискретное преобразование Фурье, реализуемое с помощью быстрого преобразования Фурье (БПФ). В результате получается комплекснозначный спектр $F_{b,c}(u,v)$. 

\begin{equation}
    F_{b,c}(u, v) = \sum_{h=0}^{H-1} \sum_{w=0}^{W-1} x_{b,c}(h, w) e^{-2\pi i \left( \frac{uh}{H} + \frac{vw}{W} \right)}
\end{equation}
где $x(h,w)$ — значение признака в точке с координатами $h$ (высота) и $w$ (ширина), а $F(u,v)$ — комплексный коэффициент на частоте с индексами $u$ и $v$.
В результате для каждой карты признаков получается комплексный массив того же пространственного размера, где действительная и мнимая части кодируют амплитуду и фазу соответствующих частотных компонент. Фаза несёт информацию о пространственном расположении структур, однако в данной работе анализируется распределение энергии по низким и высоким частотам, поэтому фазовые соотношения далее не используются. На следующем этапе выполняется переход к спектру мощности.

Вычисляется спектр мощности (энергии) как квадрат модуля:

\begin{equation}
    S_{b,c}(u, v) = |F_{b,c}(u, v)|^2
\end{equation}

где $b$ — индекс изображения в батче, $c$ — индекс канала, а $(u,v)$ — координаты в частотной плоскости. Энергетический спектр показывает, какой вклад вносит каждая пространственная частота в общую энергию сигнала, и инвариантен к фазе, что удобно для последующего усреднения.

На этом этапе усреднение по изображениям не выполняется. Если усреднить спектры по батчу сразу, то далее невозможно будет оценить разброс спектральных характеристик между отдельными изображениями.

\begin{equation}
    \bar{S}_{b}(u, v) = \frac{1}{C} \sum_{c=1}^{C} S_{b,c}(u, v)
\end{equation}

где $C$ — число каналов на рассматриваемом слое или блоке. В результате для каждого изображения получается собственная двумерная спектральная карта $\bar{S}_{b}(u,v)$, которая агрегирует частотный вклад всех каналов данного представления. Важно отметить, что усреднение спектров мощности — это не то же самое, что усреднение самих карт признаков с последующим вычислением спектра. Это важно, поскольку на одних изображениях могут преобладать горизонтальные структуры, а на других — вертикальные; усреднение карт признаков привело бы к взаимному уничтожению таких разнородных признаков. 

Полученный после предыдущих шагов двумерный спектр мощности $\bar{S}_{b}(u,v)$ содержит информацию не только о величине пространственной частоты, но и о направлении частотного вектора. Однако для целей данного исследования требуется описать, как меняется распределение энергии по мере увеличения пространственной частоты независимо от того, горизонтальная это частота, вертикальная или диагональная. Поэтому угловая зависимость далее не рассматривается, а двумерный спектр переводится в радиальное представление.

Пусть $A_j$ — множество точек частотной плоскости, попадающих в $j$-е радиальное кольцо:
\begin{equation}
A_j = \{(u,v): j \leq r(u,v) < j+1\},
\end{equation}
где радиус частотной точки после сдвига нулевой частоты в центр спектра определяется как
\begin{equation}
r(u,v)=\sqrt{(u-u_0)^2 + (v-v_0)^2},
\end{equation}
а $(u_0,v_0)$ — координаты центра спектра после операции \texttt{fftshift}. Точки с близким расстоянием до центра спектра объединяются в одно кольцо, после чего двумерная частотная карта заменяется одномерным профилем по радиусу.

Поскольку размеры карт признаков и частотных плоскостей для разных слоёв модели могут различаться, прямое сравнение радиальных профилей в абсолютных единицах неудобно. Чтобы обеспечить сопоставимость спектральных характеристик слоёв с разным пространственным разрешением, вводится нормированный частотный радиус.

Максимальный радиус задаётся как
\begin{equation}
R_{\max} = \frac{\min(H,W)}{2}.
\end{equation}
Такой выбор соответствует анализу частот внутри круга Найквиста по горизонтальному и вертикальному направлениям. Нормированный частотный радиус для $j$-го кольца определяется как
\begin{equation}
\rho_j = \frac{r_j}{R_{\max}},
\end{equation}
где $r_j$ — радиус, соответствующий $j$-му радиальному кольцу.

В работе используются два близких, но не одинаковых радиальных представления. Для визуализации спектральных кривых применяется средняя мощность в кольце:
\begin{equation}
M_{b,j} = \frac{1}{|A_j|}\sum_{(u,v)\in A_j}\bar{S}_{b}(u,v).
\end{equation}
Такой вариант удобен для графиков, потому что он показывает средний уровень спектральной мощности на данном радиусе и не делает дальние кольца визуально доминирующими только из-за большего числа точек.

Для расчёта численных метрик используется суммарная энергия кольца:
\begin{equation}
E_{b,j} = \sum_{(u,v)\in A_j}\bar{S}_{b}(u,v).
\end{equation}
Этот вариант корректен для вычисления доли энергии, потому что кольца разного радиуса содержат разное число частотных точек. Если использовать только среднее значение по кольцу, то дальние кольца будут недовзвешены, и такая величина уже не будет строгой долей энергии.

Однако сравнение таких функций для разных слоёв непосредственно по шкале $E(\rho)$ часто оказывается затруднительным, поскольку абсолютная энергия может существенно различаться в зависимости от глубины. Например, в глубоких слоях ResNet количество каналов растёт, а пространственное разрешение падает. Суммарная энергия карт признаков может изменяться на порядки из-за масштабирования, вносимого нормализацией, функциями активации и остаточными связями.

Для корректного сравнения слоёв с разной абсолютной энергией производится нормализация:
\begin{equation}
\tilde{E}_{b,j} = \frac{E_{b,j}}{\sum_j E_{b,j}}.
\end{equation}
Величина $\tilde{E}_{b,j}$ интерпретируется как доля энергии, приходящаяся на $j$-й радиальный диапазон для $b$-го изображения. В результате получается одномерная функция $\tilde{E}_{b}(\rho)$, определённая на интервале $[0,1]$ и представляющая собой нормированное распределение энергии по радиальным диапазонам. Значение $\rho=0$ соответствует постоянной составляющей; после вычитания среднего она обнулена, поэтому значение $E(0)$ близко к нулю. Значение $\rho=1$ соответствует границе круга Найквиста, то есть максимальной частоте, которую может нести карта признаков данного разрешения.

Сумма в знаменателе является полной энергией всех частотных компонент после удаления DC-компоненты. Нормализация превращает спектр в дискретное распределение долей энергии по частотным диапазонам. Это распределение далее используется при расчёте спектрального центроида и доли низкочастотной энергии.

У данного подхода есть важный нюанс: при нормализации теряется информация об общей энергии слоя. В работе нормализованные спектры используются для качественного анализа формы распределения: смещения пика в сторону низких частот, островершинности или, наоборот, более плоского профиля. Для количественной оценки дополнительно вводятся метрики, которые вычисляются на основе $\tilde E_b(\rho)$: спектральный центроид и доля энергии в низких и высоких частотах.

Слои в ResNet и Swin имеют разное пространственное разрешение. Чтобы корректно сравнивать их спектры, применяется приведение частотной оси к единому масштабу. При сравнении кривых спектральной энергии для моделей одного семейства за эталон принимается разрешение первого слоя, и для всех остальных слоёв частотная шкала пересчитывается таким образом, чтобы физические частоты (в циклах на изображение) были сопоставимы. Это позволяет строить графики, на которых видно, как меняется распределение энергии по одним и тем же пространственным частотам по мере углубления в сеть. При входном изображении $224\times224$ карта признаков \texttt{layer1} имеет размер $56\times56$, поэтому внутри круга Найквиста максимальный радиус равен $R_{\max}^{ref}=28$, а число радиальных отсчётов составляет $R_{\max}^{ref}+1=29$.

\begin{equation}
    \rho_{mapped} = \rho \cdot \frac{L_{layer}}{L_{ref}},
\end{equation}

где $L_{layer}$ — число радиальных отсчётов для рассматриваемого слоя, а $L_{ref}$ — число радиальных отсчётов для эталонного уровня \texttt{layer1}. 

После такого пересчёта поздние слои ResNet и Swin занимают меньший диапазон по оси частот, поскольку их пространственное разрешение ниже. Это не ошибка, а ожидаемое следствие дискретизации: у карты меньшего размера доступен меньший диапазон пространственных частот относительно исходной эталонной шкалы. Именно поэтому на графиках кривые глубоких стадий заканчиваются раньше, чем кривые начальных слоёв. Для каждого изображения метрики считаются на собственной радиальной сетке слоя после пересчёта радиуса к общему масштабу. 
 
\section{Метрики}

\textbf{Спектральный центроид}

Спектральный центроид представляет собой средневзвешенную частоту, где весом выступает нормализованная энергия. Для дискретной радиальной спектральной кривой $E(\rho_k)$ центроид вычисляется как:

\begin{equation}
    C = \frac{\sum \rho_k \cdot \tilde E(\rho_k)}{\sum \tilde E(\rho_k)}
\end{equation}

Центроид принимает значения на отрезке [0,1] и показывает, где сосредоточена основная энергия. Низкое значение центроида свидетельствует о том, что энергия представления сильнее сосредоточена в низкочастотной области - представление сглажено и содержит крупномасштабные структуры; чем выше значение, тем больше вклад средних и высоких пространственных частот - в представлении много мелких деталей, текстур и резкие границы.

Эта метрика обладает инвариантностью к масштабу, что позволяет сравнивать слои с разным пространственным разрешением и разные архитектуры. Кроме того, она чувствительна к изменениям формы спектра: даже небольшие смещения распределения энергии заметно меняют значение центроида. Поэтому центроид также удобен для вычисления скорости спектрального сжатия.

\textbf{Доля энергии в низких и высоких частотах}

Помимо среднего, можно численно измерить, как энергия распределена в хвостах распределения. Для этого выбирается пороговое значение $\tau$ (например, $\tau = 0.15$ для низких частот). Доля энергии в низких частотах:
\begin{equation}
    LowFrac = \sum_{k:\rho_k<0.15} \tilde E(\rho_k)
\end{equation}


Эта метрика характеризует, какая доля спектральной энергии сосредоточена в заданном частотном диапазоне. Рост LowFrac с глубиной соответствует большей концентрации энергии в области малых частотных радиусов. 

\textbf{Скорость «сжатия» спектра по глубине}

Также в качестве метрики можно использовать зависимость центроида от индекса слоя или относительной глубины. Далее эта характеристика называется спектральным сжатием. Спектральное сжатие – это процесс последовательного подавления высоких частот по мере углубления сети. Для его количественной оценки можно использовать линейную или экспоненциальную модель. 

Для компактного описания динамики используется экспоненциальная модель: она позволяет отразить ситуацию, когда центроид быстрее меняется на ранних уровнях и затем замедляет изменение к концу сети.

\begin{equation}
    C_{i}=C_{0} \cdot e^{-\lambda \cdot i} \label{eq:slope}
\end{equation}

Параметры $С_0$ и $\lambda$ подбираются методом нелинейных наименьших квадратов по значениям центроида на разных уровнях. В таблицах приводится коэффициент $\lambda$ – экспоненциальная скорость затухания. Эта величина безразмерна. За один слой центроид уменьшается в $e^{-\lambda} $ раз. Если $\lambda>0$, центроид убывает по глубине; если $\lambda \approx 0$, выраженного экспоненциального сжатия не наблюдается; если $\lambda<0$, центроид в среднем растёт.

Другие метрики, такие как энтропия спектра или пиковая частота, не использовались, так как они менее непосредственно отвечают поставленной цели данного исследования. Энтропия спектра характеризует «размытость», но не указывает, в какую сторону смещён спектр, поэтому она более подходит для исследования устойчивости. Пиковая частота игнорирует форму распределения, не даёт информации о ширине спектра и более чувствительна к шуму.

Для оценки разброса спектральные характеристики рассчитывались отдельно для каждого изображения выборки. После этого для каждого слоя и каждой модели вычислялись среднее значение и стандартное отклонение по изображениям. На графиках стандартное отклонение отображает разброс радиального спектрального профиля между изображениями, а в таблицах — разброс соответствующей метрики по выборке. Исключение составляет коэффициент спектрального сжатия $\lambda$, поскольку он получен из средних центроидов, а его погрешность оценивалась приближённо.

\section{Экспериментальная постановка}

Эксперименты проводились на предобученных моделях компьютерного зрения трёх семейств: ResNet, Vision Transformer и Swin Transformer. Для реализации использовались библиотеки \texttt{torch 2.10.0+cpu}, \texttt{torchvision 0.25.0+cpu}, \texttt{timm 1.0.21}, \texttt{numpy 2.1.3}. Модели семейства ResNet загружались из \texttt{torchvision.models}, а модели ViT и Swin — из библиотеки \texttt{timm}. 

Для ResNet использовались модели: \\ \texttt{resnet18}~\cite{hf_resnet18_tv_in1k}, \texttt{resnet34}~\cite{hf_resnet34_tv_in1k}, \\ \texttt{resnet50}~\cite{hf_resnet50_tv2_in1k}, \texttt{resnet101}~\cite{hf_resnet101_tv2_in1k}. \\ Во всех случаях веса задавались через параметр \texttt{weights=DEFAULT}. Для \texttt{resnet18} и \texttt{resnet34} это соответствует весам \texttt{IMAGENET1K-V1}; для \texttt{resnet50} и \texttt{resnet101} — весам \texttt{IMAGENET1K-V2} в используемой версии torchvision. 

Для ViT рассматривались модели: \\ \texttt{vit-tiny-patch16-224}~\cite{hf_vit_tiny_patch16_224_augreg_in21k_ft_in1k}, \texttt{vit-small-patch16-224}~\cite{hf_vit_small_patch16_224_augreg_in21k_ft_in1k}, \\ \texttt{vit-base-patch16-224}~\cite{hf_vit_base_patch16_224_augreg_in21k_ft_in1k}, \texttt{vit-large-patch16-224}~\cite{hf_vit_large_patch16_224_augreg_in21k_ft_in1k}. 

Для Swin использовались: \\ \texttt{swin-tiny-patch4-window7-224}~\cite{hf_swin_tiny_patch4_window7_224_ms_in1k}, \texttt{swin-small-patch4-window7-224}~\cite{hf_swin_small_patch4_window7_224_ms_in1k}, \\ \texttt{swin-base-patch4-window7-224}~\cite{hf_swin_base_patch4_window7_224_ms_in1k}, \texttt{swin-large-patch4-window7-224}~\cite{hf_swin_large_patch4_window7_224_ms_in22k_ft_in1k}. \\ Для моделей из \texttt{timm} использовался режим \texttt{pretrained=True}, то есть стандартные предобученные веса, заданные в конфигурации соответствующей модели.

Для ResNet анализировались выходы четырёх основных стадий: \texttt{layer1}, \texttt{layer2}, \texttt{layer3} и \texttt{layer4}. Для ViT и Swin спектры сохранялись для всех выбранных блоков, однако для основных графиков использовались пять уровней по глубине: начальный блок, блок на глубине $n/4$, средний блок, блок на глубине $3n/4$ и последний блок, где $n$ — число блоков модели. Для итоговых таблиц метрик использовались только три уровня: начальный, средний и глубокий. Это сделано для того, чтобы таблицы оставались сопоставимыми между архитектурами.

В качестве входных данных использовалась подвыборка ImageNet 10K \cite{priyerana_imagenet10k}, содержащая по 10 изображений для каждого из 1000 классов ImageNet. Все изображения проходили одинаковую предобработку: изменение размера, центральное кадрирование до $224\times224$, преобразование в тензор и нормализация по стандартным статистикам ImageNet:
$ \mu = (0.485, 0.456, 0.406), \quad
\sigma = (0.229, 0.224, 0.225).$
Размер батча во всех экспериментах составлял 32.

Для оценки разброса результатов спектральные характеристики рассчитывались не только в виде одного среднего значения. Для каждого изображения строилось радиальное распределение энергии выбранного представления, после чего для каждой метрики вычислялись среднее значение и стандартное отклонение по выборке:
\begin{equation}
\overline{M} = \frac{1}{N}\sum_{i=1}^{N} M_i,
\end{equation}
\begin{equation}
\sigma_M = \sqrt{\frac{1}{N-1}\sum_{i=1}^{N}(M_i-\overline{M})^2},
\end{equation}
где $M_i$ — значение метрики для $i$-го изображения, а $N$ — число изображений, участвующих в расчёте. В таблицах значения приводятся в формате $\overline{M}\pm\sigma_M$. На графиках погрешность отображается через область стандартного отклонения вокруг средней спектральной кривой либо через интервалы погрешности, если график строится по отдельным значениям метрик.

Отдельно проводились контрольные эксперименты для ResNet и Swin. Их задача состояла в том, чтобы отделить изменение спектрального профиля представлений от тривиального эффекта уменьшения пространственного разрешения. В ResNet пространственный размер карт признаков уменьшается при переходе между стадиями, а в Swin аналогичную роль играет операция слияния патчей. Поэтому наблюдаемое смещение спектра к низким частотам нельзя автоматически считать свойством только обученных признаков.

Первый контрольный эксперимент представляет собой базовый контрольный вариант на основе снижения разрешения. В нём раннее представление модели искусственно приводится к разрешению более глубокого уровня с помощью операции уменьшения разрешения, после чего для него рассчитываются те же спектральные метрики. Для ResNet сравниваются реальные \texttt{layer2}, \texttt{layer3}, \texttt{layer4} и представление \texttt{layer1}, искусственно приведённое к тем же пространственным размерам. Для Swin аналогично сравниваются реальные поздние стадии и раннее представление, приведённое к их разрешению. Если реальные глубокие представления близки к такой базовой модели, то низкочастотное смещение в основном объясняется снижением разрешения. Если же реальные представления дополнительно смещены к низким частотам, значит, часть эффекта связана уже с изменением самих признаков. Искусственное уменьшение разрешения в контрольных экспериментах выполнялось с помощью адаптивного среднего пулинга. В реализации использовалась функция \\ \texttt{torch.nn.functional.adaptive\_avg\_pool2d}, где выходной размер задавался равным пространственному размеру соответствующего глубокого слоя. Для ResNet исходным представлением служил layer1, который приводился к размерам layer2, layer3 и layer4. Для Swin аналогично использовалось начальное представление первой стадии, приводимое к разрешениям последующих стадий. Такой оператор выбран как простая необучаемая базовая процедура понижения разрешения. При этом сам средний пулинг подавляет часть высокочастотных компонент, поэтому контроль следует интерпретировать как базовую модель влияния уменьшения разрешения, а не как полное воспроизведение внутренних операций ResNet или Swin.

Второй контрольный эксперимент — внутристадийное сравнение. В нём сравниваются блоки, имеющие одинаковое пространственное разрешение. Для ResNet такой контроль выполняется внутри одной стадии, например внутри \texttt{layer3} в ResNet50. Для Swin сравниваются блоки внутри одного этапа. Если спектральные метрики меняются при фиксированном размере карты или решётки, то этот эффект уже нельзя объяснить одним только уменьшением разрешения.

Для ViT отдельный контроль понижения разрешения между блоками не требуется в том же виде, что для ResNet и Swin, поскольку стандартный ViT сохраняет одну и ту же решётку патчевых токенов по всей глубине модели. Поэтому для ViT дополнительно фиксируется, что пространственный размер токенной решётки остаётся постоянным, а наблюдаемые изменения спектра между блоками не связаны с межстадийным уменьшением разрешения.
